\paragraph{Conjugate-coordinate class: the Jacobi commutator of the plane-wave field with the horizontal fields, and the flux-resummed next-to-leading symbol (conditional; bounded no-go; supersession of the preceding Landau model).}
Keep the setting of the preceding paragraphs: Poisson limit, leaf \(\mathbb{CP}^{1,2}\), \(T^a=M^{a4}\), \(W_a\) the Hamiltonian field of \(T^a\), \(Y=q_\mu\{y^\mu,\cdot\}\), \(\phi=q\cdot y\), \(\kappa_a=W_a(\phi)\), \(A^a=\epsilon\alpha^a\cos\phi\), the nilpotent (leading) symbol with centre flux \([\omega_a,\omega_b]=\mathrm iF_{ab}\), \(\omega_a=-\mathrm iW_a\), and the term \(+M_{\mathrm{str}}\operatorname{Tr}[(H_X-\mathrm i0)^{-1}]\). The Lie--Poisson relations give exactly \(\kappa^a=\sigma M^{45}q^a\) and \([W_a,Y]=\sigma q_a\operatorname{Ham}(y^4)\), \(y^4=M^{45}\), with \(\sigma=\pm1\) the sign convention of the bracket (the second is the Jacobi identity; both checked at random orbit points to \(10^{-14}\)). In the frame \((W_a,Z_k)\), \(Z_k\) a basis of the vertical space, \(Y=y^cW_c+y^kZ_k\) and the field \(V=\operatorname{Ham}(y^4)\) has no vertical component (below \(10^{-14}\) at six sampled points), so \([W_a,Y]\) has no component along the centre. The flux term \(y^cF_{ac}\) is not small and is not parallel to \(\kappa_a\) (the residual of the parallel fit is \(0.7\)--\(1.0\) at the sampled points and directions of the centre variable).

\emph{Consequence for the model of the preceding paragraph.} There the plane-wave coupling was \(\zeta_Y=y^c\omega_c+c\) with a constant charge \(c\), for which \([\omega_a,\zeta_Y]=\mathrm iy^cF_{ac}\ne0\); this commutator is inconsistent with the Jacobi identity at the nilpotent order with frozen coefficients. That model is not invariant under \(\alpha\to\alpha+\lambda\kappa\), the leaf form of \(\delta A^a=\{\Lambda,T^a\}\), at \(F\ne0\) and \(y\ne0\) (relative change \(0.07\)--\(0.16\) of the second-order response in a test, \(10^{-14}\) at \(F=0\)), whereas the Fin values quoted there are values of that model. On the leaf \(Y\) commutes with \(\omega_a\) and with \(\phi\) (\(Y(\phi)=0\)); at the nilpotent order it is a right-invariant field of the nilpotent group, and its diagonal value at the base point is \(P=y^c\omega_c+c\), \(c=y^k\zeta_k\), which stands at the seam of the trace, \(\operatorname{Tr}[\text{chain}\cdot P^2]\), the chain being the ordered product of the left-invariant factors. Because \(P^2\) separates the first and the last propagator, the time-translation reduction \(\int_0^1\mathrm du\,(1-u)\langle V(su)V(0)\rangle\) is not available and the second order is the double Duhamel integral \(-s^2\int_{0<u_2<u_1<1}\langle V(su_1)V(su_2)P^2\rangle\).

\emph{Checks.} In a two-dimensional Euclidean Landau problem in Landau gauge, where the magnetic translation \(K=y_2p_2\) is an exact c-number, this prescription reproduces the second-order density obtained by exact diagonalisation (relative difference \(2\cdot10^{-4}\), real); the prescription with \(P^2\) at the end of a merged propagator does not (complex, \(1.5\%\) off in the real part). The two-dimensional test validates the prescription for a central \(K\); the identification of the leaf field \(Y\) with the seam insertion \(P\) is conditional on the nilpotent order and the frozen first jet. In four dimensions the central model is invariant under \(\alpha\to\alpha+\lambda\kappa\) to \(10^{-14}\) at \(F\ne0\), \(y\ne0\), and equals the model of the preceding paragraph at \(F\to0\) (\(10^{-8}\)). On the leaf, with \(F_{ab}(\zeta)=-S_{ab}{}^k\zeta_k\), \(\zeta\) the centre variable, \(f_e/f_m=\ell/|y^4|\) at three sampled directions of \(\zeta\) at each of five orbit points, the quadratic form of the Fierz--Pauli plus dilaton coefficient \(A\) of the next-to-leading-symbol paragraph equals the closed form in \((\alpha,\kappa,y^c)\) used there (ratio \(1\) to six digits, five orbit points).

\emph{Result in the class.} (i) The coefficient of \(\mathrm ds/s\) is unchanged: \(c=A/16\pi^2\), with \(T_2(s)\,s/c\to1.0000\) as \(s\to0\) at the flux of the leaf, in both the central and the preceding model. The Jacobi commutator therefore does not change, cancel or add to the Einstein--Hilbert-form coefficient \(A\) of the logarithm. (ii) At \(F\ne0\) the integrand decays as \(\mathrm e^{-|s|(f_e\cos\vartheta+f_m\sin\vartheta)}\) on \(s=|s|\mathrm e^{-\mathrm i\vartheta}\), and the finite part \(\mathrm{Fin}=\lim[\int_{s_0}^\infty T_2\,\mathrm ds-c\ln(1/s_0)]\) is a finite complex number for each centre variable. Split by the content of \(P\) into \(y^cy^d\) (tetrad), \(c^2\) (charge) and cross terms, the three parts are separately invariant under \(\alpha\to\alpha+\lambda\kappa\) (\(10^{-7}\) or better). (iii) The tetrad part of Fin, as a quadratic form in \(\alpha\), is not proportional to the form \(A\): the residual of the best proportional fit, in real and imaginary parts, is \(0.54\) and \(0.85\) (centre direction \(0.8\)), \(0.08\) and \(0.48\) (direction \(2.0\)) at one orbit point, and \(0.004\) and \(0.000\) at a second point where \(f_e/f_m\simeq42\); the quotient by \(\kappa\) has eigenvalues of both signs in every case (as has \(A\)). (iv) The form \(A\) is not Lorentz covariant, as shown for the preceding paragraph (it depends on \(y^\mu\) and on the fibre point), and the flux does not alter the coefficient that carries this dependence. In the sampled per-centre-variable forms the tetrad part of the finite remainder is neither proportional to the form \(A\) nor sign-definite; the sign and covariance of the full finite remainder (tetrad, charge and cross parts) and of its window average over the centre variable were not tested, so no finite, covariant, sign-definite Einstein--Hilbert coefficient is obtained here.

\emph{Size of the dropped terms.} The horizontal part of \([W_a,Y]\) (the weight-one term \(-q^av^cW_c\)), the horizontal parts of \([W_a,W_b]\), the gradients of the flux and of \(y^c,y^k\), and the window are not included. The ratio of the norm of the horizontal part of \([W_a,Y]\) to that of \(y^cS_{ac}{}^k\) is \(0.01\)--\(0.5\) at six orbit points (per unit centre variable), so these terms are small only for \(|\zeta|\) large compared with that ratio and for small \(\omega\); this is the regime of the nilpotent approximation, not a bound on the leaf.

This is a statement for the Poisson limit, one loop, the nilpotent (leading) symbol resummed in the flux with the Jacobi-consistent central plane-wave field, the Lorentzian \(\mathrm i\epsilon\) proper time, the \(M_{\mathrm{str}}\) term, and numerical evaluation at sampled orbit points and centre directions. It does not include the vector-mass term with the Zeeman-dressed resolvent at order \(\epsilon^2\) (whose leading growth was found above and is not shown to be altered), the weight-one terms, the window average over the centre variable, the sign and covariance of the full finite remainder, the spectral definition of the vector sector, or non-perturbative masses, and it identifies no induced form with the emergent metric of Chapter~\ref{ch:vierbein-emergence}, makes no Wick-rotation claim, and gives no value or sign of \(G_{\mathrm{eff}}\). This changes no ledger label; Problem~F remains open.
